% abstract: plain scientific abstract (taste pass 2026-10-07). Plain words; definitions live in the Introduction and Table I.
\begin{abstract}
Swarms of language-model agents now work together for months. The people who study them and the people who run them both need a description far smaller than the logs that still predicts what the swarm will do. We ask which physics models give such a description for one swarm, the AI Village: 46 agents, 51 shared goals and 18 months of logs. We matched 17 candidate models against 132 pre-registered hypotheses, each with its prediction, its no-effect baseline and its kill condition fixed before it touched real data, and each required to rule out the shared inputs that make agents move together without reading each other. Two models hold up well, two hold in part, and together they give one picture. Each agent is a kinetic spin that changes state only at its own model calls. Agents influence each other only through messages that a call reads, mostly messages that name the reader, and that influence is weak: one message causes about 0.13 more, so echoes die out, and the swings of total talk match this gain with no free parameter. Strong external fields set most of the state: the daily timetable, the goal text, the chat rooms and each agent's own style. A message gives a fast kick that fades in about 7 calls, on top of a slow return to the agent's own resting point over about 100. Beyond one shared mode, few patterns of joint movement rise above chance. We tested and mostly dropped the superagent and information-dynamics lines of inquiry. Seven predictions that were written down before their tests, among them the fluctuation rule and the moment of influence, are compared with their outcomes in plain terms. From the picture we derive 16 ranked quantities for operators to watch. All results are exploratory; tests on reserved data have confirmed almost none of them.
\end{abstract}
